\documentclass[11pt]{article}
\usepackage[letterpaper,margin=0.6in,top=0.55in,bottom=0.6in]{geometry}
\usepackage[T1]{fontenc}
\usepackage{lmodern,microtype,enumitem,xcolor,fancyhdr,array}
\renewcommand{\familydefault}{\sfdefault}
\setlist{nosep,leftmargin=1.2em,itemsep=1.5pt}
\setlength{\parindent}{0pt}\setlength{\parskip}{4pt}
\pagestyle{fancy}\fancyhf{}\renewcommand{\headrulewidth}{0pt}
\fancyfoot[L]{\footnotesize Bellingham Math Circle / Week 24 / Adult table card / Unpiloted}
\fancyfoot[R]{\footnotesize Full answers: the facilitator guide}
\newcommand{\sect}[1]{\par\medskip{\large\bfseries #1}\par\smallskip}
\begin{document}
{\footnotesize Week 24 / Three random decks / For the adult at each table}\par
{\LARGE\bfseries Three random decks: the one-page version}\par
\sect{What the activity is}
Three little decks of number cards beat each other in a circle, like rock, paper, scissors: A beats B, B beats C, and C beats A. Children play short games with the decks, then list every pair of cards that could come up to find out which deck is really ahead.

\sect{Make the decks (once, about 30 minutes; keep them)}
\begin{itemize}
\item Each table gets three paper lunch bags labelled \textbf{A}, \textbf{B} and \textbf{C}, each holding three cards made from half index cards, number written large:
\par\smallskip
\begin{tabular}{>{\bfseries}l l l l}
Bag A & 2 & 4 & 9\\ Bag B & 1 & 6 & 8\\ Bag C & 3 & 5 & 7
\end{tabular}
\item K--1 table: also draw each number's dots in a ten-frame (two rows of five), as on page~1 of the K--1 packet.
\item Spares: one more copy of each of the nine cards per table, in pencil (later problems copy or rewrite cards). Grades 2--3 table: also cards 0, 10, 11 and 12, and the letters A, B, C on sticky notes rather than on the bags (Problem~3 hides them). The room needs 58 cards.
\item K--1 table: 18 counters, no wider than 1/2 inch. Everyone: pencils and blank paper.
\end{itemize}

\sect{Launch (one minute, whole group)}
Hold up bags A and B. Without looking, draw one card from each; the bigger number wins the round for its deck. Put both cards back and shake. Ask: \emph{``How could we find every pair of cards that could come up, even ones we never draw today?''} Don't say which deck wins.

\sect{What Problem 1 is for}
\begin{itemize}
\item \textbf{K--1:} partners play three games, A against B, B against C and C against A, six rounds each, putting a counter in the winning deck's row. The games raise the question ``is one deck best?'' They can't settle it: the stronger deck wins a six-round game only about 45\% of the time. Accept any result. Problems 2--4 settle it by circling the winner of all nine card pairs.
\item \textbf{Grades 2--3 and 4--5:} the ``pick first or pick second'' game, ten rounds. Each row of the table is one game: write both deck letters and their rounds won, like the shaded example row (``A, 4 wins'' and ``C, 6 wins''). The idea to reach: whoever picks second can always take the deck that beats the first pick. Ten rounds often won't show it; Problem~2 proves it by counting pairs.
\end{itemize}

\sect{The surprise}
Each deck beats the next one on 5 of its 9 card pairs, so no deck is best, even though every deck adds up to 15.

\sect{Key answers}
\begin{itemize}
\item \textbf{K--1.} P2--4: A beats B, B beats C, C beats A, each 5 pairs to 4. P5: no; the stronger deck can still lose six rounds. P6: only 9. P7: five swaps work, the whole row of B's 1 and the whole column of A's 9. P8: impossible, since nine pairs can't split into two equal piles. P9: impossible, since the deck holding the 1 loses at least two of its four pairs to every other deck.
\item \textbf{Grades 2--3.} P2 as above. P3: no number of rounds makes the guess certain. P4: A 1, 5, 9 / B 3, 4, 8 / C 2, 6, 7, and A 2, 3, 9 / B 1, 6, 8 / C 4, 5, 7. P5: 2 or 4. P6: doubling every card changes nothing; with the four-card decks C beats both others, so the circle breaks. P7: e.g.\ A 3, 5, 7 / B 2, 4, 12 / C 1, 9, 11. P8: impossible, as K--1 P9.
\item \textbf{Grades 4--5.} P2: 5/9 each way round; equal totals of 15 don't decide it. P3: exactly the three arrangements above. P4: e.g.\ A 1, 4, 4 / B 3, 3, 3 / C 2, 2, 5. P6: a two-card deck wins when its small card beats their small card and its big card beats their big card. P7: three cards per deck is the fewest. P8: impossible.
\end{itemize}
\end{document}
